\documentclass[11pt,a4paper]{article}
\usepackage[top=2.2cm,bottom=2.2cm,left=2.3cm,right=2.3cm]{geometry}
\usepackage[colorlinks=true,urlcolor=blue,linkcolor=blue,citecolor=blue]{hyperref}
\usepackage{amsmath,amssymb}
\usepackage{booktabs}
\usepackage{xcolor}
\usepackage{enumitem}
\usepackage{titlesec}
\usepackage[round,authoryear]{natbib}

\definecolor{rcol}{rgb}{0.16,0.16,0.16}
\definecolor{acol}{rgb}{0.05,0.25,0.55}

\newcommand{\eps}{\varepsilon}
\newcommand{\F}{\mathcal{F}}

% R: reviewer point (quoted/paraphrased).  A: our response.
\newenvironment{point}[1]{%
  \par\medskip\noindent\textbf{\color{rcol}#1}\par\nopagebreak\smallskip\begingroup
  \color{rcol}\itshape}{\endgroup\par}
\newenvironment{reply}{%
  \par\smallskip\noindent\textbf{\color{acol}Response.}\ \begingroup}{\endgroup\par}

\titlespacing*{\section}{0pt}{2.2ex plus 1ex minus .2ex}{1.2ex plus .2ex}
\titlespacing*{\subsection}{0pt}{1.8ex plus 1ex minus .2ex}{1ex plus .2ex}

\begin{document}

\begin{center}
{\Large\bfseries Response to Reviewers}\\[0.6em]
{\large VOUCH: Anytime-Valid Certificates for Machine Unlearning via
Paired-Canary Hypothesis-Betting}\\[0.4em]
\emph{Machine Learning} --- revised submission
\end{center}

\bigskip

\noindent We thank both reviewers for reports that were unusually careful and, in two
places, mathematically decisive. Reviewer~1's observation that Definition~1's marginal
null does not match the condition our supermartingale proof actually uses is correct,
and acting on it has changed the paper's central theorem rather than merely its
wording. Reviewer~1's correction of the Kullback--Leibler expansion is likewise
correct. Reviewer~2's insistence that every end-to-end result sat at a permissive
tolerance, that the declared score class was never stress-tested from outside, that the
$\eps=0.02$ entries in Table~3 were censoring artefacts, and that we had asserted the
canaries ``do not stand out'' without testing it, were each right, and in the last case
the measurement we now report contradicts what we had written. We have made every
change requested. Where a request was not fully satisfiable within the compute
available we say so explicitly rather than substituting a weaker experiment.

\medskip
\noindent\textbf{Summary of the substantive changes.}
\begin{enumerate}[leftmargin=1.4em,itemsep=1pt,topsep=2pt]
\item \textbf{The certificate's null is restated and its theorem replaced.} We now
certify the \emph{realised advantage on the committed cohort} and bet against a
without-replacement boundary recentred at every step. The resulting guarantee
(Theorem~2) is exact under arbitrary dependence across pairs, templates and repetition
strata --- no independence, no exchangeability across pairs, no homogeneity. A new
Proposition~3 prices the transfer to the super-population advantage under a named
assumption. Both targets are reported side by side throughout.
\item \textbf{New simulation showing the old null was genuinely not controlled}
(\S6.1, Tables~1--2): under model-level over-dispersion the previous fixed-boundary
process issues under the marginal null in $38.9\%$ of runs, an eightfold breach of
$\alpha$; the corrected process holds at $\le 0.020$ in every regime. Heterogeneity is
also measured on the real runs.
\item \textbf{Four valid sequential comparators added} (\S6.4, Table~5): exact
group-sequential $\alpha$-spending (Pocock, O'Brien--Fleming), a Jeffreys-mixture
e-process, a Robbins normal-mixture e-process, and a fixed-$n$ test at a pre-committed
size.
\item \textbf{Tolerance sweep to $\eps = 0.05$ on real models} (Table~9), with the
cohort/super-population distinction made explicit; guidance on choosing $\eps$ added.
\item \textbf{Mathematical corrections}: $\mathrm{KL} = \eps^2/2 + O(\eps^4)$, so the
cohort rule is $2\log(1/\alpha)/\eps^2$; the score sign convention fixed to the
token-normalised \emph{log-likelihood}; per-score versus simultaneous confidence-sequence
coverage stated; Proposition~7's one-directionality stated immediately after it.
\item \textbf{Statistical reporting}: Clopper--Pearson intervals on every verdict count,
across-seed dispersion on every mean column, censoring rates and Kaplan--Meier medians
in the power table, seed counts on the model axis, and the ``stable across the axis''
characterisation corrected.
\item \textbf{Three new empirical sections}: head-to-head against forget-quality and
min-$k\%$ on identical runs (\S6.10); attacks from \emph{outside} the declared class
(\S6.11); canary detectability, with the negative result and the entropy-dilution
design rule that follows from it (\S6.12).
\item \textbf{SimNPO added} as a subject; \textbf{benchmark forget/retain and
general-capability readings} added; \textbf{reproducibility appendix} added, including
an unambiguous definition of ``retraining''.
\item \textbf{Scope made consistent} in abstract, introduction, experiments and
conclusion; legal claims sourced to primary instruments; the $99.6\%$ figure attributed
to our own earlier design; promotional phrasing removed; the six papers requested are
discussed, not merely cited.
\end{enumerate}

\clearpage
\section*{Reviewer 1}

\subsection*{1. The null in Definition~1 does not match the condition used in Theorem~2}

\begin{point}{R1.1}
Definition~1 defines $p^{(s)} = \Pr[Z_i^{(s)}=1]$, a marginal probability, whereas the
proof assumes $\mathbb{E}[Z_i^{(s)}\mid\mathcal{F}_{i-1}] \ge p_0$ at every step, which
is stronger. The signs may not be independent because all canaries pass through the
same jointly trained and unlearned model, and may differ across templates and
repetition strata. Please either redefine the certificate using an appropriate
conditional-mean null and justify why the setting satisfies it, or provide a result
valid under the dependence structure considered here. It would also help to state which
sources of randomness define $p^{(s)}$. Simulations with dependent or heterogeneous
signs would strengthen this part.
\end{point}

\begin{reply}
This is correct, it is the most consequential point in either report, and we have taken
the second of the two options offered --- a result valid under the dependence structure
--- because we agree the first would have been a patch. Four changes follow.

\emph{(a) Which randomness defines the null (\S4.3, \S5.1).} The randomness is the
\emph{inclusion coins}, and we now make the filtration carry that literally. The
manifest is committed as a Merkle tree with one leaf per pair, and leaf $\pi(t)$ is
opened at step $t$, so the coin $b_{\pi(t)}$ is genuinely unrevealed at time $t-1$.
Theorem~1 is restated against this filtration and now reads: under exact unlearning,
$Z^{(s)}_{\pi(t)}\mid\mathcal{F}_{t-1}\sim\mathrm{Bern}(\tfrac12)$ \emph{exactly},
\textbf{whether or not the signs are independent across pairs}. The proof never asks
for independence; it asks only that the coin governing the pair opened now still be
fair given everything opened before, which under exact unlearning it is. The revocation
arm therefore needs no repair --- the reviewer's concern does not bite there, and we now
say so explicitly.

\emph{(b) The certificate is retargeted (Definition~2, Theorem~2).} The composite null
``some score still leaks'' is where dependence does bite. Once $M_u$ and the manifest
exist, the cohort's sign vector is \emph{fixed}; the only randomness left to the audit
is the reveal order. We therefore certify
\[
\Delta^{(s)} \;=\; \frac{2}{m}\sum_{i=1}^m z_i^{(s)} - 1,
\]
the \emph{realised cohort advantage} --- what the declared score actually achieves
against the committed cohort, which is what an auditor of \emph{this} model wants
bounded. Sequential reveal of a committed uniform permutation is simple random sampling
without replacement, so we bet against the recentred boundary
$p_0(t) = (m p_0 - \sum_{u<t} Z_{\pi(u)})/(m-t+1)$, the mean the unrevealed remainder
would need for the null to hold. Theorem~2 shows this is a nonnegative supermartingale
\emph{conditionally on the sign vector}, so its coverage holds under \textbf{arbitrary}
dependence across pairs, arbitrary heterogeneity across templates and repetition
strata, and arbitrary non-identical distribution. No assumption whatsoever is spent on
how the sign vector arose.

\emph{(c) The super-population claim is available but priced (Proposition~3).} Under a
named \emph{bounded coin influence} assumption --- flipping one coin changes at most
$\kappa$ signs, with $\kappa=1$ automatic under exact unlearning --- McDiarmid gives
$|\Delta^{(s)} - \bar\Delta^{(s)}| \le \kappa\sqrt{2\log(2/\alpha')/m}$. At $m=512$,
$\alpha'=0.05$, $\kappa=1$ that is $0.12$. So the algorithm-level statement costs a
tolerance inflation we now state, and this is exactly why our tight-tolerance
end-to-end results are cohort certificates (see R2.W1).

\emph{(d) Simulations with dependent and heterogeneous signs (\S6.1, Tables~1--2).}
We built the study you asked for and it vindicates the concern. Sign streams with
\emph{marginal} mean exactly $p_0=0.55$ under four regimes, $2{,}000$ seeds, peeking
after every pair:

\begin{center}\small
\begin{tabular}{lccc}
\toprule
regime & $\mathrm{sd}(\hat\Delta)$ & fixed boundary (v1) & without replacement (ours)\\
\midrule
independent                & 0.031 & 0.034 & 0.014\\
heterogeneous by stratum   & 0.030 & 0.026 & 0.020\\
model-level over-dispersion& 0.296 & \textbf{0.389} & 0.001\\
shared model-level latent   & 0.500 & \textbf{0.481} & 0.000\\
\bottomrule
\end{tabular}
\end{center}

\noindent Independence and stratum heterogeneity alone leave the old process valid.
A shared model-level random effect --- one training run leaks, another does not --- breaks
it eight- to tenfold. The corrected process controls the realised-cohort error at
$\le 0.020$ everywhere. Table~2 verifies Theorem~2 directly on \emph{fixed}
adversarially structured sign vectors at the least favourable feasible point of the
cohort null: error never exceeds $0.035$ against nominal $0.05$, for every structure,
both tolerances, both cohort sizes.

Finally, the dependence is real on our own runs: a chi-square test of sign-rate
homogeneity across repetition strata rejects in $23$ of the $27$ un-unlearned runs,
where there is leakage to be heterogeneous about, and in $16\%$ of runs overall.

We note in \S5.2 one honest consequence of the retargeting: a cohort audited to
exhaustion resolves deterministically, so the betting machinery buys early stopping,
validity under continuous monitoring, and the alarm arm, rather than buying the cohort
statement itself. We prefer to state that than to leave the reader to discover it.
\end{reply}

\subsection*{2. Scope of the certificate}

\begin{point}{R1.2}
VOUCH certifies the advantage of a declared score class on the planted canary
population, which is not the same as proving all organic forget records were removed.
Some statements, such as ``certifies genuine unlearning,'' sound stronger than the
guarantee. Make the scope consistent throughout. Provide more direct evidence that the
canary results track forgetting on organic records. Relatedly, the commitment does not
prevent a provider who knows the canaries from treating them differently; state the
honest-but-verifiable assumption whenever the certification claim is summarised.
\end{point}

\begin{reply}
Agreed on all three counts, and acted on.

\emph{Consistency.} ``Certifies genuine unlearning'' is gone. The abstract now closes
with ``The certificate concerns a declared score class on a committed cohort under an
honest-but-verifiable provider, and we measure what each qualification costs.'' The
introduction has a dedicated paragraph, \emph{What the certificate claims, and what it
does not}, listing the four qualifications with the measured cost of each. \S5.4 is
retitled \emph{Scope, stated precisely}. The conclusion enumerates the same four with
their costs. We also checked every remaining occurrence of ``certif*'' in the
manuscript for an unqualified claim.

\emph{Evidence that canaries track organic forgetting.} We have added the comparison
you suggest, in the strongest form our data supports (\S6.10, Table~10): on identical
models and seeds we compute a TOFU-style forget-quality analogue (two-sample KS between
the unlearned and retrained models' canary-score distributions) and a MUSE-style
min-$k\%$ leakage analogue (paired AUC relative to the retrained reference), and report
them beside VOUCH's verdicts. Three disagreements emerge, and we report them in both
directions. The forget-quality analogue is systematically \emph{pessimistic},
failing genuine unlearning methods whose realised advantage is within noise of zero,
because a KS test detects distributional differences far below any tolerance one would
certify --- it answers ``are these identical'' when the question is ``is the residual
advantage below $\eps$''. PrivLeak reads $-1.2\%$ on the MUSE/Phi-1.5 GradDiff seed
that VOUCH \emph{revokes}, which is the concrete instance of a claim the previous
version made only rhetorically. And min-$k\%$ inspected after every pair flags leakage
on genuinely clean retrained models in several cells. We are also even-handed about
the limits: these are analogues computed on canary cohorts, not the benchmarks' own
metrics on their own splits, and we say so.

We have additionally strengthened the dose--response reporting, since it is the only
calibration of the canary-to-organic extrapolation we can honestly offer. Mean score
gaps rise monotonically with training-time repetition ($+0.42$, $+0.60$, $+1.14$,
$+1.85$ nats at $r = 1,2,4,8$) and so do realised advantages ($+0.28$, $+0.32$,
$+0.54$, $+0.70$). We now also state the limit this implies: the $r=1$ stratum, the one
most like an organic record, carries the \emph{weakest} signal, so the extrapolation is
weakest exactly where it matters most. We do not claim a transfer theorem, because we
do not have one.

\emph{The threat model.} The honest-but-verifiable assumption is now stated in \S4.1,
in the introduction's scope paragraph, in \S5.4, and in the conclusion --- i.e.\ wherever
the certification claim is summarised. It also acquired teeth we had not expected: see
our response to R2.8, where a measurement shows a zero-knowledge perplexity filter
recovers the \emph{entire} planted cohort. We retract the manuscript's previous claim
that the canaries ``do not stand out'' and now explain both why detection does not by
itself defeat the audit (both twins are equally conspicuous, so locating the cohort
reveals nothing about which twin the coin admitted) and why closing the remaining gap
requires cryptographic attestation, which we do not supply.
\end{reply}

\subsection*{3. Mathematical and notational inconsistencies}

\begin{point}{R1.3a --- sign convention}
Eq.~(2) defines $D_i = s(M_u, c_{\mathrm{in}}) - s(M_u, c_{\mathrm{ghost}})$, but
Fig.~5 appears to use the reverse sign: NLL values $0.99/9.03$ are reported as
$D = +8.0$, which under Eq.~(2) should have the opposite sign if $s$ is NLL. Use one
convention throughout theory, figures, tables and implementation.
\end{point}

\begin{reply}
You identified a real inconsistency, and we can now say precisely where it was. The
released implementation's score is \texttt{s\_loss = mean(logprobs)}, i.e.\ the
token-normalised \emph{log-likelihood} --- the \emph{negative} of the NLL. Under that
definition $D = (-0.99) - (-9.03) = +8.04$, so Eq.~(2), Fig.~5, the tables and the code
were all mutually consistent; what was wrong was the \emph{prose} in \S4.2, which
described the score as ``the token-normalised negative log-likelihood''. That single
word inverted the stated convention.

We have fixed the prose rather than the figure, and pinned the convention down so it
cannot drift again. \S4.4 now lists every member of $\F$ with the orientation stated
explicitly (``each oriented so that a \emph{higher} score means \emph{more
memorised}''), defines $s_{\mathrm{loss}}$ as the token-normalised log-likelihood, and
adds a sentence recording that an earlier draft said otherwise. Equation~(2) is now
followed by the reading rule: a positive $D$ means the trained twin is the
better-remembered one, which is the direction of leakage. Figure~5's caption and body
now restate the convention at the point of use. We verified the $+8.0$ figure against
the released pipeline output.
\end{reply}

\begin{point}{R1.3b --- the small-$\eps$ expansion}
$\mathrm{KL}(\mathrm{Bern}(1/2)\|\mathrm{Bern}(1/2+\eps/2)) = -\tfrac12\log(1-\eps^2)
= \eps^2/2 + O(\eps^4)$, so the leading approximation for certification time should be
about $2\log(1/\alpha)/\eps^2$, not $8\log(1/\alpha)/\eps^2$. Interestingly the
corrected expression is consistent with Table~3 and with ``about 600 pairs for
$\eps=0.1$, $\alpha=0.05$''. Please correct the theorem and recheck related
calculations.
\end{point}

\begin{reply}
You are exactly right, including in your inference about Table~3. Theorem~4 now carries
the derivation in full:
\[
\mathrm{KL}\bigl(\mathrm{Bern}(\tfrac12)\,\|\,\mathrm{Bern}(p_0)\bigr)
= -\tfrac12\log\bigl(4p_0(1-p_0)\bigr) = -\tfrac12\log(1-\eps^2)
= \tfrac{\eps^2}{2} + O(\eps^4),
\]
so $\mathbb{E}[\tau^*] \approx 2\log(1/\alpha)/\eps^2$. We rechecked every dependent
number. The tables were computed in code from the \emph{exact} divergence
(\texttt{kl = 0.5*log(0.5/p0) + 0.5*log(0.5/(1-p0))}), so no numerical column changes:
at $\eps=0.1$, $\alpha=0.05$ the exact bound is $596$ and $2\log(1/\alpha)/\eps^2$ gives
$599$, against the erroneous rule's $2{,}397$. Only the stated rule of thumb was wrong,
by a factor of four. Theorem~4 now says so in a sentence, so a reader of the previous
version is not left guessing. The cohort-sizing prose is updated ($600$ pairs at
$\eps=0.1$, $2{,}400$ at $0.05$, $15{,}000$ at $0.02$), as is the cost discussion, which
previously quoted the wrong scaling.
\end{reply}

\begin{point}{R1.3c --- per-score or simultaneous confidence sequences?}
Clarify whether the confidence-sequence guarantee is per score or simultaneous over the
whole score class $\F$.
\end{point}

\begin{reply}
Added as Remark~5, distinguishing three statements. Each score's sequence covers
\emph{its own} advantage with uniform-in-time probability $1-\alpha$: the guarantee is
per score. The reported $\max_{s\in\F} U_t^{(s)}$ \emph{is} nonetheless a valid
level-$\alpha$ anytime upper bound on $\sup_{s\in\F}\Delta^{(s)}$ with no multiplicity
correction, because once the cohort is fixed the maximising index of the population
quantities is a fixed, non-random index, and only that one score's sequence is needed.
Simultaneous two-sided coverage of \emph{all} $|\F|$ advantages does cost a union bound
and holds at $|\F|\alpha$; we now report intervals at that corrected level wherever the
whole class is displayed.
\end{reply}

\subsection*{4. A fairer optional-stopping comparison}

\begin{point}{R1.4}
The fixed-sample binomial test inspected after every pair is a useful illustration but
not a strong baseline. Please add at least one valid sequential baseline --- alpha-spending,
group-sequential, or another Bernoulli confidence-sequence/e-process method --- under the
same observation stream and query budget. A standard fixed-$n$ comparison at a
precommitted sample size would also help show the practical cost of anytime validity.
\end{point}

\begin{reply}
Agreed; the peeking binomial was a demonstration, not a competitor, and we have added
all four comparators you name (\S6.4, Table~5), on the same sign stream and the same
query budget:
\begin{itemize}[leftmargin=1.4em,itemsep=0pt,topsep=2pt]
\item exact group-sequential $\alpha$-spending with Pocock and O'Brien--Fleming
spending functions, with boundaries calibrated by \emph{dynamic programming over the
exact binomial lattice} at the least favourable null rather than by a normal
approximation, so the baseline is as strong as it can be made;
\item a truncated Beta$(\tfrac12,\tfrac12)$-mixture (Jeffreys) e-process;
\item a Robbins normal-mixture e-process;
\item a fixed-$n$ binomial test at a pre-committed sample size.
\end{itemize}

Every procedure holds its level. Against exact unlearning at $\eps=0.1$ with a budget
of $1{,}536$ pairs: VOUCH (cohort) certifies $100\%$ of runs using a mean of $523$
queries; Pocock $97.1\%$ using $645$; O'Brien--Fleming $98.5\%$ using $844$; the mixture
e-processes $77$--$78\%$; the fixed-$n$ test $98.9\%$ but spending all $1{,}536$ queries
on every run, since it cannot stop early.

The honest summary, which we now give in the text, is that anytime validity is close to
free here: against the strongest valid sequential comparator VOUCH gives up nothing in
power and saves about a fifth of the queries, and against fixed-$n$ it gives up one
point of power for a threefold reduction in verifier cost. The table also isolates how
much of the advantage comes from our retargeting rather than from betting: on the
super-population target, same budget, power falls from $1.000$ to $0.882$ and mean
queries rise from $523$ to $747$ --- the same price Proposition~3 charges in tolerance.
We think separating those two effects is more useful than a single headline number.
\end{reply}

\subsection*{5. Stronger evidence for the broader conclusions}

\begin{point}{R1.5a --- three seeds and counts like 16/18}
Most benchmark results use three seeds, while conclusions are drawn from counts such as
16/18 and 17/18. Report uncertainty where possible and use more cautious wording when
the number of runs is small.
\end{point}

\begin{reply}
Done, and the wording is now built around the intervals rather than the counts. A new
Table~10 gives every pooled verdict count with a Clopper--Pearson $95\%$ interval:
$16/18 \to [0.65, 0.99]$, $18/18 \to [0.81, 1.00]$, $17/18 \to [0.73, 1.00]$. Mean
columns carry across-seed standard deviations. Every claim in \S6.9 is now stated as
``$16$ of $18$ runs, $95\%$ CI $[0.65,0.99]$'' rather than as a bare fraction, and the
model-axis discussion is rewritten (see R2.W6) to stop treating single-seed rows as
evidence of stability.
\end{reply}

\begin{point}{R1.5b --- SimNPO and another recent method}
SimNPO is discussed in related work but not evaluated. Adding SimNPO and, if possible,
another recent utility-preserving method would make the comparison more representative.
\end{point}

\begin{reply}
SimNPO \citep{fan2025simnpo} is implemented and evaluated. We added it as a
reference-free objective,
$L = (2/\beta)\,\mathrm{softplus}(-\beta(\mathrm{nll}_\theta - \gamma))$ on the
length-normalised forget NLL with $\beta=0.1$, $\gamma=0$ --- NPO's objective with the
frozen reference replaced by a constant margin --- and ran it end to end on TOFU/GPT-2
under identical seeds, canary manifests and budgets as the existing subjects, so it is
directly comparable. It certifies on all three seeds and achieves the widest
forget/retain separation of any subject that remains fluent (forget NLL $6.20$ against
retain $2.62$, versus NPO's $5.50$ against $2.60$), at a capability NLL of $4.17$
against retraining's $3.23$. It is reported in the new metrics table
(Table~12) and named in \S6.9 as the best-behaved of the approximate methods on the
utility axis.

On a second method we must be candid about the boundary of what this revision achieved.
We prioritised SimNPO because you named it and because it is the closest competitor to
a subject we already had. Adding RMU \citep{li2024wmdp} or task-vector negation would
require a further round of fine-tuning across all six benchmark/model cells, and we
judged the marginal value lower than the four new \emph{measurement} studies
(\S6.1, \S6.4, \S6.10, \S6.11, \S6.12) that speak to the framework's validity rather
than to its coverage of the method space. We name this as a limitation rather than
implying the method list is complete.
\end{reply}

\begin{point}{R1.5c --- held-out NLL is not sufficient for utility}
Held-out NLL identifies severe collapse but does not establish general utility.
Consider reporting standard TOFU/MUSE forgetting and retention measures together with
at least one general capability benchmark.
\end{point}

\begin{reply}
Agreed, and the added metrics changed what we can say. The end-to-end runner now
records, behind an \texttt{--extra-metrics} flag and enabled for the revision-round
runs: mean NLL on the benchmark's own \emph{forget} and \emph{retain} splits;
greedy-decode ROUGE-L recall against the reference continuation on both splits, the
TOFU/MUSE reporting convention; and a general-knowledge \texttt{capability\_nll} on a
fixed probe set --- TOFU \texttt{world\_facts}, or the MUSE holdout news stream.
Results are in the new Table~12, on TOFU/GPT-2 over three seeds.

They sharpen the paper's central utility finding rather than softening it. Gradient
ascent reaches forget NLL $94.3$, retain NLL $93.7$ and ROUGE-L recall of exactly
$0.00$ on \emph{both} splits: it has not forgotten the forget set, it has stopped
producing language. More usefully, the metrics reveal something held-out NLL alone had
hidden: GradDiff's retain NLL is $6.8$ but its capability NLL is $19.4$, so its damage
is nearly threefold larger \emph{off} the retain distribution than on it. A
retain-only utility reading --- which is what we previously had --- would have missed
that entirely, which is a direct vindication of your point.

Two honest limitations. First, \texttt{world\_facts} is a modest general-knowledge
probe, not a suite such as MMLU; we say so in Appendix~A and name a broader suite as
the next step, preferring a small clearly described probe to implied coverage. Second,
we discovered a confound while running this and report it rather than printing the
flattering number: for TOFU the P1 relearn corpus \emph{is} the capability probe set,
so a post-P1 capability reading measures training on the probe. It reads $1.46$,
better than the retrained reference, which is an artefact. That cell is marked
unavailable in Table~12 with the reason given.
\end{reply}

\begin{point}{R1.5d --- systematic R-VOUCH}
The R-VOUCH experiments would be more convincing if relearning, quantisation and
jailbreak probes were evaluated more systematically across the main unlearning methods
under clearly specified attack budgets.
\end{point}

\begin{reply}
The budgets are now specified exactly, which was the more serious of the two gaps:
Appendix~A states P1 $=$ $100$ optimiser steps on a public corpus that never contains a
canary, at batch $4$ and learning rate $1\times10^{-4}$; P2 $=$ symmetric per-channel
round-to-nearest 4-bit quantisation of all linear and embedding weights; P3 $=$ three
fixed adversarial wrappers applied at scoring time only. These budgets are identical
across every subject and model, which was true before but undocumented.

On systematic coverage across methods we have made partial progress and prefer to
report it as such. The probes remain anchored on NPO in the benchmark tier, where
P1 and P3 run on all six model/benchmark cells and all three seeds; P2 runs on the
GPT-2 and TinyGPT tiers, and is skipped for the multi-billion-parameter models because
the shared-frozen-base multi-adapter design would have to materialise merged weights to
quantise them. Extending all three probes to all five unlearning methods multiplies the
benchmark tier by roughly three in training cost, which was beyond this revision's
budget. We have removed any implication that the probe study is exhaustive and state
the coverage explicitly.
\end{reply}

\subsection*{6. Related work}

\begin{point}{R1.6}
The P1 relearning probe is closely related to Hu et al., ICLR 2025. The evaluation
discussion would benefit from Wang et al., ICLR 2025 and Yuan et al., ICLR 2025, and
Li et al., ICLR 2026 is particularly relevant to apparent forgetting and misleading
automatic metrics. Please discuss these works in relation to VOUCH rather than only
adding citations.
\end{point}

\begin{reply}
All four are now discussed substantively, in a restructured \S3.3 titled
\emph{Evaluating unlearning: a literature of negative results}, which we organised into
three strands so that each work's bearing on VOUCH is explicit rather than implied.

\citet{hu2025jogging} is placed as the direct antecedent of our P1 probe: their
demonstration that benign relearning on unrelated public text restores both hazardous
knowledge and verbatim passages is precisely the obfuscation-versus-removal distinction
P1 operationalises, and it is why we treat a post-relearning verdict as part of the
certificate rather than an optional extra. We connect it to our own measurement that
the probe moves NPO from issued to undetermined on TOFU/GPT-2 and TOFU/Pythia and
\emph{revokes} on one MUSE/Phi-1.5 seed.

\citet{wang2025towards} contributes two findings we now act on: forget-quality metrics
track surface behaviour and collapse under red-teaming, and the reported
unlearning/retention trade-off is not a Pareto frontier, so method rankings are
hyperparameter artefacts. We now state explicitly that this is why VOUCH does
\emph{not} rank unlearning methods against one another but certifies each against a
declared tolerance with a paired utility reading.

\citet{yuan2025closer} shows ROUGE-style matching hides degenerate models. We connect
this to our own Figure~5, in which an NPO-unlearned model emits \texttt{LLLLLL}\ldots{}
on canary prompts while held-out loss is untouched --- an instance of exactly their
failure mode, caught here not by a fluency metric but by two-sided certification, which
treats scoring conspicuously \emph{below} the ghost twin as leakage.

\citet{li2026beliefs} required the most careful treatment because it bears on the
\emph{limits} of what we certify, and we have written it that way rather than
defensively. Their squeezing effect --- probability mass pushed off the target and
redistributed onto semantically equivalent rephrasings --- means mass that has migrated
to a paraphrase is, by construction, mass our default score class does not see, since
our scores read the likelihood of the literal secret span. We state this plainly, note
that the framework's response is structural (the class $\F$ is declared, the
certificate is explicitly a statement about $\F$, and any named score can be added at
the cost of power and no loss of validity since Theorem~2 needs no multiplicity
correction), point to our new out-of-class experiments (\S6.11) as a first step, and
identify a paraphrase-aware score in their spirit as the single most valuable extension
of this work. It is also cited in the conclusion's limitations.

We additionally added, on your and Reviewer~2's prompting, a new \S3.2
\emph{Measuring memorisation and forgetting} placing VOUCH in the secret-sharer lineage
\citep{carlini2019secret} and the forgetting-measurement lineage
\citep{jagielski2023measuring}. The comparison there is one we should have drawn in the
first submission: exposure answers ``how memorised is this string,'' whereas the ghost
twin answers ``could this have happened by chance,'' and the latter is what a
certificate needs and what \citet{zhang2024mia} show requires controlled randomisation.
\end{reply}

\subsection*{7. Reproducibility and the cost comparison}

\begin{point}{R1.7a --- what ``retraining from scratch'' means}
Please define exactly what ``retraining from scratch'' means in the LoRA experiments.
If this means training a fresh adapter from the same frozen base model without the
forget data, the terminology should make this clear.
\end{point}

\begin{reply}
Your reading is exactly right, and the terminology was misleading. Appendix~A now opens
with the definition: initialise a \emph{fresh} adapter on the same frozen base and
train it on $D_{\mathrm{keep}}$ only --- neither organic forget data nor any canary ---
with the identical optimiser, schedule and step count as the original fine-tune. It is
exact unlearning \emph{within the adapter paradigm}: the adapter that saw the canaries
is discarded, and the frozen base never saw a canary at any point. It is not retraining
of pretraining, which we never perform. The subject is renamed
\textbf{``retrain (fresh adapter)''} in every table.

We also answer, in the same appendix, R2.11's related question about whether
adapter-level unlearning is representative of the deployment setting: the audit is
indifferent, since it queries $M_u$ as a black box, but a rank-$16$ adapter has fewer
degrees of freedom in which to hide memorised content than a full fine-tune, so our
subjects are certified in a regime that if anything favours them. We flag full-parameter
runs as the natural next scaling step.
\end{reply}

\begin{point}{R1.7b --- reproducibility settings}
Please provide the key settings: exact model checkpoints, LoRA target modules, learning
rates, training and unlearning steps, query counts, canary fractions, ordering
procedure, tie-breaking rules, and stopping criteria.
\end{point}

\begin{reply}
Appendix~A now gives all of them: the eight exact model identifiers (including that
Phi-1.5 is forced to fp32 because its LoRA fine-tuning diverges to NaN in fp16); LoRA
rank $16$, $\alpha_{\mathrm{LoRA}}=32$, dropout $0$, \texttt{target\_modules =
all-linear}; the exact TOFU and MUSE splits and formatting; canary domains, repetition
strata, secret entropy, per-tier cohort sizes and the measured canary character share;
fine-tune/retrain at $600$ steps, batch $8$, AdamW $2\times10^{-4}$, block $160$, clip
$1.0$, and every unlearning stage's steps and hyperparameters individually; $Q=2$
wrappers for benchmark and model-axis runs and $Q=4$ for the synthetic tier, with the
score class per tier; the reveal order (a permutation seeded by the run seed, committed
before any query); the tie-break PRNG seed ($20260702$, shared across all runs); the
stopping criteria for both arms, including that the revocation arm deliberately
continues after a crossing so the streaming product alarm stays well defined; and the
seed counts for every simulation tier.
\end{reply}

\begin{point}{R1.7c --- the cost comparison}
Claims such as ``roughly sixty times cheaper'' or ``two orders of magnitude cheaper''
should be based on clearly matched hardware and workloads. Report the cost assumptions
more explicitly.
\end{point}

\begin{reply}
This criticism was well founded: the previous version quoted a single
``seventy-six seconds'' without naming the configuration and compared it to
alternatives on unmatched hardware. Both are now fixed and the speed-up framing is
abandoned.

Verification cost is stated as a workload: exactly $2Qm$ short forward passes and no
training --- $1{,}536$ to $2{,}560$ passes per subject at our settings. Measured
wall-clock medians are reported \emph{per tier, on the hardware that tier ran on}:
$13$\,s (TinyGPT), $63$\,s (GPT-2, $640$ pairs), $112$--$148$\,s (the 124M--1.4B
benchmark tier), $397$--$499$\,s (the 3.8B--5.1B models). The comparison to descriptive
alternatives is now a ratio of \emph{workloads on the same model and hardware}: a
retrain-based reference costs one full fine-tune, a shadow panel at least sixteen,
against VOUCH's zero; on our TOFU/GPT-2 configuration one fine-tune is $600$ steps at
batch $8$, about $4{,}800$ sequence forward-and-backward passes, so a sixteen-model
panel is roughly $77{,}000$ training passes against VOUCH's $1{,}536$ inference passes.
Appendix~A states the heterogeneous hardware explicitly and explains why we give a
ratio rather than a speed-up factor: the two workloads are not interchangeable.
\end{reply}

\subsection*{Minor comments}

\begin{point}{R1--M1 --- GDPR and EU AI Act}
Some statements about the GDPR and EU AI Act are broad. Please cite appropriate primary
sources or use more careful wording.
\end{point}

\begin{reply}
Rewritten against the primary instruments, both now cited with their Official Journal
references and CELEX identifiers \citep{gdpr2016,aiact2024}. The wording is corrected in
two ways that matter. First, Article~17 GDPR creates the erasure \emph{right} but
contains no demonstrability requirement; the duty to be able to \emph{show} compliance
lives in Articles~5(2) and 24(1), with supervisory audit powers in Article~58(1)(b), and
we now attribute it there. Second, we no longer write that the AI Act has ``audit
provisions,'' which overstates it: the Act imposes automatic event logging
(Article~12), technical documentation and ten-year retention (Articles~11, 18),
production to authorities on reasoned request (Article~21), conformity assessment
(Article~43), and gives the AI Office evaluation powers over general-purpose models
(Article~92). The paragraph now closes by noting that neither instrument prescribes a
statistical test, which is the actual gap the paper addresses.
\end{reply}

\begin{point}{R1--M2 --- promotional phrasing}
Phrases such as ``dissolves all three obstacles,'' ``forgetting by lobotomy,'' and
``what a regulator cares about'' are somewhat promotional for a journal paper.
\end{point}

\begin{reply}
All three removed, along with others we found on a sweep: ``dissolves'' $\to$
``addresses''; ``forgetting-by-lobotomy'' $\to$ ``destructive unlearning''; ``it is
verification, not optimisation, that a regulator cares about'' $\to$ ``Evaluation, not
optimisation, is the bottleneck,'' now supported by citations rather than assertion.
We also toned down ``impossible to miss,'' ``emphatically,'' ``straw man'' and
``embarrassment'' in the experiments section.
\end{reply}

\begin{point}{R1--M3 --- dense figures and tables; Table~7 seed counts}
Some figures and tables are dense and use small text. Table~7 should make it
immediately clear which models use one seed and which use three.
\end{point}

\begin{reply}
Table~7 (the model axis) now carries an explicit \textbf{seeds} column, and its
surrounding text is rewritten to distinguish the three-seed rows from the single-seed
ones instead of describing the whole axis as stable (see R2.W6). We split the former
single benchmark table into three narrower ones --- verdicts, the tolerance sweep, and
the interval-estimated counts --- which removes the worst of the density, and the new
tables are set at \texttt{\textbackslash small} rather than \texttt{\textbackslash
tiny}. The benchmark verdict table remains the densest object in the paper at
\texttt{\textbackslash tiny}; it carries six model/benchmark cells and we could not find
a narrower layout that keeps them side by side, which is the comparison the table
exists to support.
\end{reply}

\clearpage
\section*{Reviewer 2}

\noindent We are grateful for a report that engaged with the construction closely
enough to find the Table~3 censoring artefact and to press on the $\eps=0.2$ default.
Both were right.

\subsection*{W1. Every end-to-end result uses $\eps = 0.2$}

\begin{point}{R2.W1 / Point 1}
$\eps=0.2$ means an adversary can still distinguish the in-twin from its ghost about
$60\%$ of the time. Tighter tolerances are never run on an actual language model. Report
end-to-end results on at least one real model at $\eps \le 0.05$, or state explicitly in
the abstract, introduction and conclusion that all end-to-end certification is at
$\eps = 0.2$, and discuss what that tolerance means operationally. Guidance on how to
choose $\eps$ would also help.
\end{point}

\begin{reply}
We have done the first of the two, and also the discussion you asked for.

\emph{End-to-end at $\eps \le 0.05$ (Table~9).} The corrected finite-cohort target
(R1.1) is materially more powerful than the fixed-boundary one, and replaying every
saved run at $\eps \in \{0.05, 0.1, 0.2\}$ now yields real-model certificates at
$\eps = 0.05$: exact unlearning is issued on all three seeds at $\eps=0.05$ for
TOFU/Pythia-160M and TOFU/Phi-1.5, and on $13$ of $18$ benchmark runs overall. At
$\eps = 0.1$ it is issued on all $18$.

\emph{And the caveat that must travel with it.} This is a \emph{cohort} certificate. On
the super-population target no run certifies at $\eps=0.05$ at these cohort sizes, and
Theorem~4 says none should: $\eps=0.05$ needs roughly $2{,}400$ pairs and we plant at
most $640$. Table~9 therefore reports both targets side by side, and we recommend in the
text that they be read together. We think this is a more useful answer than either
overclaiming or planting $3{,}000$ canaries in a $5{,}000$-record academic corpus, where
the cohort would be $37\%$ of the training data and no longer a plausible model of
deployment.

\emph{What $\eps = 0.2$ means, said plainly.} \S6.9 now contains an explicit
subsection on choosing $\eps$, which states that $\eps=0.2$ licenses a $60\%$
distinguishing rate and $\eps=0.05$ a $52.5\%$ one, and gives three considerations that
set it: a regulatory one (if the audit must support a demonstrability claim, the
permitted advantage should not itself be reportable as a privacy loss, arguing for
$\eps \le 0.05$); a statistical one (Theorem~4 fixes the cohort at
$\approx 2\log(1/\alpha)/\eps^2$ pairs, so $\eps$ is bought with canaries --- $600$ at
$0.1$, $2{,}400$ at $0.05$, $15{,}000$ at $0.02$); and an operational one (the cohort must
stay a negligible corpus share, which makes $\eps=0.02$ affordable for a provider with
millions of records and unaffordable at benchmark scale). We now describe our own
$\eps=0.2$ as a \emph{demonstration} tolerance chosen so that $384$--$640$ pairs suffice
on academic corpora, explicitly not a recommendation, wherever the number appears.

\emph{One further consequence we surface rather than hide.} Under the corrected, more
powerful target, the two GPT-2 cells in which the un-unlearned model was previously
``undetermined'' are now \emph{certified} at $\eps=0.2$ --- their realised advantages are
$0.062$ and $0.078$, genuinely below the tolerance. So at $\eps=0.2$ a model that was
never unlearned at all can pass. That is the sharpest available illustration of your
point about the permissiveness of the default, and we now use it as such in \S6.9.
\end{reply}

\subsection*{W2. Canaries versus a user's deleted record}

\begin{point}{R2.W2 / Point 2}
The certificate concerns canaries; the motivating scenario concerns a user's deleted
record. The introduction opens with GDPR Article 17 and the abstract makes no
qualification. The $r$ strata calibrate but do not license the extrapolation, and the
$r=1$ end is where the measured gap is smallest. Either give a transfer statement under
an explicit named assumption, or reposition the contribution as a certificate about a
planted cohort. Also clarify that the guarantee is per-cohort, not per-request.
\end{point}

\begin{reply}
We have repositioned, and additionally supplied a transfer statement --- but for the
other transfer, since the two should not be conflated.

\emph{Repositioning.} The paper now says throughout that the certificate is about the
committed canary cohort. The abstract's closing sentence names all three
qualifications. The introduction's scope paragraph states that the guarantee is
per-cohort, \emph{not per-deletion-request}, because per-request certificates would need
per-user canaries --- and this now appears in the introduction, not only in \S5.4 and the
conclusion. The legal framing is rewritten (R1--M1) so that it no longer implies the
certificate discharges an Article~17 obligation for a particular data subject; it
presents the accountability duty as the reason evidence is needed at all.

\emph{The transfer statement we can prove.} Proposition~3 bridges the \emph{realised
cohort} advantage to the \emph{super-population} advantage over the coin draw, under an
explicitly named bounded-coin-influence assumption, at cost
$\kappa\sqrt{2\log(2/\alpha')/m}$. That is a genuine, named, quantified transfer.

\emph{The transfer we cannot prove, stated as such.} The step from planted canaries to
\emph{organic} records is a different matter and we do not have a theorem for it. We
say so, and we have strengthened the empirical calibration and its stated limits: the
dose--response curve is now reported with dispersion, and we state explicitly that
$r=1$, the stratum most like an organic record, carries the weakest signal, so the
extrapolation is weakest exactly where it matters most. \S6.10's head-to-head against
forget-quality and min-$k\%$ on identical runs (see R1.2) is the most direct evidence we
can offer that canary verdicts track conventional forgetting measurements, and we
present it with its limitations rather than as a substitute for the missing theorem.
\end{reply}

\subsection*{W3. Proposition 6 runs one way only}

\begin{point}{R2.W3 / Point 3}
Proposition 6 shows $(\eps_u,\delta_u)$-certified algorithms pass VOUCH. There is no
converse, so passing VOUCH implies no $(\eps_u,\delta_u)$ guarantee. Calling it a
``bridge'' and ``semantic anchor'' invites the wrong reading. State the
one-directionality immediately after the proposition.
\end{point}

\begin{reply}
Done, as Remark~8 placed immediately after the proposition, and stated at least as
strongly as you put it: ``The converse is false and we claim nothing like it: passing
VOUCH does \emph{not} imply any $(\eps_u,\delta_u)$ guarantee, does not bound the
influence of any datum outside the canary cohort, and does not constrain scores outside
$\F$.'' The remark also reframes the proposition's purpose as a consistency check on
the definition --- an algorithm with a parameter-space guarantee should not be flagged by
a behavioural audit --- and says it should not be read as a semantic equivalence in
either direction. The phrase ``semantic anchor'' is deleted; \S3.4 now calls it a
consistency check and cross-references the remark. The proposition is also retitled
``One-directional bridge from certified unlearning'' so the qualification travels with
the name.
\end{reply}

\subsection*{W4. The declared class $\F$ is never stress-tested from outside}

\begin{point}{R2.W4 / Point 4}
The certificate is a supremum over $\F$, so a stronger attack outside $\F$ could exceed
$\eps$ on a certified model. The R-VOUCH probes re-run the same score class, so they do
not test this. Run at least one attack outside $\F$ against a certified model and report
whether the residual advantage stays below $\eps$. If infeasible, say why, and soften
the claim that VOUCH bounds ``the largest residual advantage any score in the class can
extract.''
\end{point}

\begin{reply}
We ran the experiment, it produced a partial negative result, and we report it as such
(\S6.11, Table~11). Two attacks that no member of $\F$ implements were run against all
$157$ model/seed combinations certified at $\eps=0.2$:
\begin{itemize}[leftmargin=1.4em,itemsep=0pt,topsep=2pt]
\item \emph{stratum-restricted}: the repetition label $r_i$ is part of the committed
manifest, so a verifier may condition on it; restricting the audit to $r=8$
concentrates it on the most-memorised pairs. No score in $\F$ uses $r_i$.
\item \emph{cross-fitted learned combination}: a Fisher direction over the declared
scores estimated on the first half of the reveal order and applied to the second. A
trained combination of members of $\F$ is not itself a member of $\F$.
\end{itemize}
Both preserve Theorem~1 --- the first conditions only on committed per-pair data, the
second uses a predictable weight vector --- so both come with anytime-valid bounds rather
than point estimates.

\begin{center}\small
\begin{tabular}{lccccc}
\toprule
attack & runs & pairs & mean $|\hat\Delta|$ & max $|\hat\Delta|$ & still certifies\\
\midrule
declared class $\F$ (reference) & 157 & --- & 0.040 & 0.130 & ---\\
stratum-restricted ($r=8$)      & 157 & 96  & 0.069 & 0.270 & \textbf{0.968}\\
cross-fitted learned combination& 157 & 192 & 0.057 & 0.172 & 1.000\\
\bottomrule
\end{tabular}
\end{center}

The learned combination never exceeds the tolerance. The stratum-restricted attack
\emph{does}, on $5$ of $157$ certified runs ($3.2\%$), reaching a realised advantage of
$0.27$ against a certified bound of $0.2$. This does not contradict Theorem~2, which
concerns $\F$ and the whole cohort, but it does show the class restriction is
load-bearing.

Accordingly we have softened the claim exactly as you ask. The phrase ``the largest
residual advantage any score in the class can extract'' is gone; the guarantee is
written as a statement about $\F$ everywhere, including in the abstract, the
introduction's scope paragraph, \S4.4, \S5.4 and the conclusion. We also give the
constructive reading: the framework absorbs such attacks with no change to the theory,
since adding the stratum-restricted score to $\F$ makes the certificate slower and
stronger and Theorem~2 needs no multiplicity correction.

Finally, on the LiRA-style attack you suggested: we did not run it. It needs dozens of
retrained shadow references per cell and was outside the compute available for this
revision. We say this in \S6.11 and name it as the most valuable next experiment,
rather than implying the two attacks we did run are a substitute.
\end{reply}

\subsection*{W5. No comparison against an existing verification method}

\begin{point}{R2.W5 / Point 5}
Section 5.5 asserts that VOUCH's alarm catches a GradDiff failure ``that a forget-metric
evaluation would have passed,'' but that comparison is never run. Adding a column
showing what forget quality and min-$k\%$ conclude on the same models and seeds would
convert the paper's central rhetorical claim into evidence, and it is cheap to do.
\end{point}

\begin{reply}
You were right that it was cheap and right that we should have done it. \S6.10 and
Table~10 now report both metrics on the same models and the same seeds, and the specific
claim you flagged is now substantiated: on the MUSE/Phi-1.5 GradDiff seed that VOUCH
\emph{revokes}, the min-$k\%$ leakage analogue reads $-1.2\%$ --- a \emph{negative}
leakage figure on a model carrying real residual memorisation. That is the evidence the
previous version asserted.

Two further disagreements appear, and we report them even though one is unflattering to
the descriptive metrics and the other is unflattering to nobody in particular. The
forget-quality analogue is systematically pessimistic on this data, failing genuine
unlearning methods whose realised advantage is within noise of zero, because a KS test
on several hundred paired scores detects distributional differences far smaller than any
tolerance one would certify. And the min-$k\%$ statistic inspected after every pair ---
the way an auditor facing a stream of deletions would naturally use it --- flags leakage
on genuinely clean retrained models in several cells, which is optional stopping doing
what \S6.2 predicts.

We are also explicit about the comparison's limits: these are analogues computed on
canary cohorts, not the benchmarks' own metrics on their own splits, and a metric
designed for one target should not be judged too harshly on another. The narrow point
stands: on identical data the descriptive and certified readings disagree in both
directions.
\end{reply}

\subsection*{W6. Thin statistical reporting}

\begin{point}{R2.W6 / Point 6}
Three seeds per cell, bare I/R/U patterns, no interval estimates, no dispersion on the
mean columns. Counts such as ``sixteen of eighteen'' carry wide binomial uncertainty at
$n=18$. Table 7 appears to use a single seed for the larger models and shows a markedly
weaker picture, yet the text describes the axis as ``stable''. The honest reading is
that the framework returns undetermined at the cohort sizes used for the models that
matter most.
\end{point}

\begin{reply}
Every part of this is addressed, and we accept the characterisation in the last
sentence.

\emph{Intervals and dispersion.} New Table~10 gives every pooled verdict count with a
Clopper--Pearson $95\%$ interval ($16/18 \to [0.65,0.99]$; $18/18 \to [0.81,1.00]$;
$17/18 \to [0.73,1.00]$). The GPT-2 end-to-end table gains an across-seed standard
deviation column on $\bar D$. Counts in the prose are now always given with their
interval.

\emph{Table 7 and the ``stable'' claim.} The table now carries an explicit
\textbf{seeds} column making the single-seed rows unmistakable, and the surrounding
text is rewritten. It no longer says the axis is stable. It says: detection is stable
--- every un-unlearned model is revoked at every scale, on every seed run; certification
is \emph{not} uniformly stable --- at the $384$-pair cohorts used for the largest models
the super-population target returns undetermined for retraining and NPO on Qwen3-0.6B
and Qwen3-4B, exactly as Theorem~4 predicts at that cohort size; and one seed per cell
is too thin to distinguish a systematic effect from noise. We also state that those
cells do issue under the cohort target, so the reader can see which claim is
cohort-size-limited and which is not.
\end{reply}

\subsection*{W7. Table 3's starred entries are selection-biased}

\begin{point}{R2.W7 / Point 7}
Starred medians are taken over runs that issued within a 20,000-pair horizon, so both
starred entries sit below the information-theoretic bound, which cannot be right and is
an artefact of censoring. As printed, the row reads as though VOUCH beats the KL limit.
Report the censoring rate and either use a survival estimate or mark them as lower
bounds.
\end{point}

\begin{reply}
This was a genuine error and your diagnosis of it is exactly correct; an entry below
the information limit should have been a red flag to us. The power table now reports the
censoring rate and a Kaplan--Meier median, and the artefact disappears:

\begin{center}\small
\begin{tabular}{lcccc}
\toprule
$\eps$ & median $\mid$ issued & Kaplan--Meier median & censored & $\log(1/\alpha)/\mathrm{KL}$\\
\midrule
0.02 & 11{,}189 & 18{,}377 & \textbf{0.45} & 14{,}976\\
0.05 & 2{,}998  & 2{,}998  & 0.00 & 2{,}394\\
0.10 & 686      & 685      & 0.00 & 596\\
0.20 & 166      & 165      & 0.00 & 147\\
\bottomrule
\end{tabular}
\end{center}

At $\eps=0.02$, $45\%$ of runs are censored and the corrected median is $18{,}377$,
comfortably \emph{above} the limit of $14{,}976$; conditioning on the fast half had
produced $11{,}189$. At $\alpha=0.01$ censoring reaches $72\%$ and the survival curve
never reaches one half, so no median exists and we now print ``---'' rather than a
number. The caption states the censoring convention, and \S6.5 explains the artefact
explicitly so a reader of the earlier version can see what changed.

One presentational note. The table also has a finite-cohort column whose entries can
legitimately sit below the KL bound, because the cohort claim resolves deterministically
once the cohort is exhausted and so is capped by $m$ rather than governed by a
super-population information limit. We flag this in both the caption and the text so
that the column is not misread the way the starred entries were.
\end{reply}

\subsection*{Point 8. Are the canaries detectable?}

\begin{point}{R2.8}
The canaries are templated with fictitious entities and 40 bits of entropy. The paper
asserts they ``do not stand out'' but never tests this. Please report a detection
experiment --- for example, whether a perplexity filter or a small classifier can separate
canaries from organic corpus records --- since a provider who can identify the committed
cohort defeats the audit, and the honest-but-verifiable threat model is doing a lot of
work here.
\end{point}

\begin{reply}
We ran the experiment and it contradicts what we had written. We retract the claim.

A perplexity filter under the \emph{pristine base model} --- one forward pass per record,
no knowledge of the template --- separates the deployed canaries from the TOFU corpus with
AUC $1.000$ and recovers the \emph{entire} cohort at a $1\%$ false-positive rate
(\S6.12, Table~12). A trained character-$n$-gram classifier does the same.

There is an information-theoretic reason, and stating it turns the negative result into
a design rule. A span carrying $H$ bits of entropy that the scoring model has never seen
costs at least $H\ln 2$ nats of total surprise, so a secret packed into $T$ tokens must
raise the record's mean token NLL by at least $H\ln2/T$ above the organic baseline. Our
deployed template spends $51.7$ bits over $25$ tokens: a predicted excess of $1.45$ nats
per token against an organic standard deviation of $0.42$. Detection was not a failure
of craft but an arithmetic certainty.

\begin{center}\small
\begin{tabular}{lccccc}
\toprule
template & bits & tokens & predicted excess & ppl.\ AUC & recall @ 1\% FPR\\
\midrule
deployed (alphanumeric) & 52 & 25 & 1.45 & 1.000 & 1.00\\
word-composed           & 42 & 35 & 0.84 & 1.000 & 1.00\\
entropy-diluted         & 68 & 88 & 0.54 & 0.974 & \textbf{0.21}\\
\bottomrule
\end{tabular}
\end{center}

The rule identifies the lever as $T$, not the vocabulary --- re-spending the budget on
natural words changes nothing --- and acting on it works: an entropy-diluted template
carrying \emph{more} entropy ($68$ bits) over $88$ tokens cuts the filter's recall from
$1.00$ to $0.21$. Both templates are in the released generator. We keep two caveats
visible: a \emph{trained} detector still recovers every cohort, because the template's
fixed phrasing is a signature --- though a provider who can fit one already knows the
template, a far stronger position than the perplexity filter needs and one the threat
model already excludes; and dilution presumably costs power, since a longer span
dilutes the memorisation signal along with the entropy, which we have not yet measured
end to end and therefore do not present as a drop-in replacement.

On your framing that a provider who identifies the cohort ``defeats the audit'', we
would refine it slightly, because the distinction turns out to be the crux of the threat
model. Both twins in a pair are drawn from the same template distribution and are
therefore \emph{equally} conspicuous, so a provider who locates the cohort learns which
records are canaries but nothing whatever about which twin the coin admitted: the null
of Theorem~1 survives complete disclosure of the manifest's \emph{contents}. What it
does not survive is a provider who, having located the cohort, treats it differently
from organic data. That is special-casing; it is what the honest-but-verifiable
assumption excludes; and it is precisely the gap cryptographic attestation exists to
close. We now state that assumption wherever the certification claim is summarised, and
say in the conclusion that we now know it is easier to violate than we had supposed.
\end{reply}

\subsection*{Point 9. Query count $Q$ and the cost figure}

\begin{point}{R2.9}
Report the number of query prompts $Q$ and how they were chosen. $Q$ affects both power
and the 76-second cost figure and is currently unspecified. State which configuration
produced the 76\,s number, and give cost as a function of $\eps$ and model size.
\end{point}

\begin{reply}
$Q$ is now stated wherever it matters. $Q = 2$ for every benchmark and model-axis run
(the identity wrapper and ``According to the archives, \ldots'') and $Q = 4$ for the
GPT-2 synthetic tier, which also carries the base-calibrated \texttt{ratio} score;
Appendix~A lists the wrappers and the per-tier score class. The wrappers were chosen as
light paraphrases of the canary prefix that preserve the secret span's position, and
$Q$ was set to $2$ in the benchmark tier for cost reasons, which we now state rather
than leave implicit.

On the cost figure: the previous ``seventy-six seconds'' is withdrawn rather than
re-attributed, because it did not correspond to a configuration we could identify
unambiguously in the logs, and quoting it would compound the original error. Cost is
now given as a closed form plus per-tier measurements. The workload is exactly $2Qm$
forward passes and no training. As a function of the two variables you name: $m$ scales
as $2\log(1/\alpha)/\eps^2$ by Theorem~4, so verifier cost scales as
$4Q\log(1/\alpha)/\eps^2$ forward passes, and the per-pass cost scales with model size.
Measured medians, per tier and on the hardware that tier ran on, are $13$\,s (TinyGPT,
0.9M), $63$\,s (GPT-2, $640$ pairs), $112$--$148$\,s (124M--1.4B benchmark tier) and
$397$--$499$\,s (3.8B--5.1B). Appendix~A records the heterogeneous hardware and explains
why we report workload ratios rather than speed-up factors.
\end{reply}

\subsection*{Point 10. The tie rate}

\begin{point}{R2.10}
State the tie rate for $D_i = 0$ empirically. With 4-bit quantised models and degenerate
NPO outputs, near-ties may be systematic rather than negligible, and the committed-coin
fix deserves a measured justification.
\end{point}

\begin{reply}
Measured across all $89{,}088$ scored pairs in the study, and the answer is more
interesting than we expected, in the direction you anticipated. For
$s_{\mathrm{loss}}$, exact ties occur in \emph{zero} pairs, at every architecture from
0.9M to 5.1B parameters, including the degenerate-output NPO runs. For
$s_{\mathrm{mink}}$ there is one exception and it is exactly where you predicted one:
$2.7\%$ of pairs ($42$ of $1{,}536$) tie exactly on Gemma-4-E2B, because the min-$k\%$
statistic over a short span takes few distinct values under that model's quantised
inference path.

So the committed tie-break coin is a formality for the loss score and is
load-bearing for min-$k\%$ on that model --- which is the measured justification you
asked for, and we report both numbers in \S4.4 and \S6.9 rather than only the reassuring
one. We also report the near-tie rate at $|D| < 10^{-6}$ in the released re-analysis
JSON.
\end{reply}

\subsection*{Point 11. LoRA adapters and the deployment setting}

\begin{point}{R2.11}
Clarify that all fine-tuning and unlearning is applied to LoRA adapters of rank 16, and
discuss whether adapter-level unlearning is representative of the deployment setting the
certificate is meant to cover.
\end{point}

\begin{reply}
Stated in the experiments preamble, in Appendix~A, and in the definition of ``retrain
(fresh adapter)'' (R1.7a). On representativeness we give a direct two-part answer
rather than a disclaimer. The \emph{audit} is indifferent: it queries $M_u$ as a black
box and never inspects parameters, so nothing in Theorem~1 or Theorem~2 changes if the
provider unlearns by full-parameter gradient steps. What the adapter setting changes is
the \emph{difficulty of the problem the subjects face}: a rank-$16$ adapter has far
fewer degrees of freedom in which to hide memorised content than a full fine-tune, so
our subjects are certified in a regime that if anything favours them, and a provider
doing full-parameter unlearning should expect verdicts at least as demanding. We flag
this as a limitation of the empirical study rather than of the framework, and name
full-parameter runs as the natural next scaling step.
\end{reply}

\subsection*{Point 12. Framing the 99.6\% figure}

\begin{point}{R2.12}
The 99.6\% figure is a comparison against the authors' own earlier design under an
adversarially constructed composite null, not a refutation of any published method. The
abstract currently reads as the latter.
\end{point}

\begin{reply}
Correct, and corrected. The abstract now reads: ``we show that our own earlier
adaptive-mixture design is anytime-invalid, falsely certifying $99.6\%$ of the time
under an adversarially constructed composite null.'' Both attributions --- ours, and
adversarial --- are now attached to the figure at every occurrence: in the abstract, in
the introduction's list of corrections, in Remark~4 (``which is our own earlier design,
not a published method''), in \S6.3 (``It is a trap we fell into and climbed out of,
reported so that others need not''), and in the caption of the soundness table.
\end{reply}

\subsection*{Minor comments}

\begin{point}{R2--m1 --- secret sharer and forgetting measurement}
The canary methodology descends from the secret-sharer line, and the
forgetting-measurement literature is closely related; neither is cited. Please add
Carlini et al.\ (2019) and Jagielski et al.\ (2023) and situate the ghost-twin design
relative to them.
\end{point}

\begin{reply}
Both added, with a new \S3.2 devoted to situating the design. \citet{carlini2019secret}
is credited with the planted-canary device and the exposure metric, and we state the one
change that matters for inference: exposure answers ``how memorised is this string,''
whereas randomising which twin enters training makes the measurement carry a
\emph{null}, and the ghost twin answers ``could this have happened by chance'' --- the
question a certificate must answer and the one \citet{zhang2024mia} show requires
controlled randomisation. \citet{jagielski2023measuring} is credited with measuring
forgetting via the decay of attack success, and we note that their two findings --- that
models do forget, but that non-convex models can retain indefinitely and that
deterministic training exhibits no forgetting at all --- are the premise of ours:
incidental forgetting is real but uncontrolled, and so cannot substitute for a
certificate. \citet{barbulescu2024each} is connected to the same discussion as the
per-sequence memorisation-strength result our repetition strata traverse.
\end{reply}

\begin{point}{R2--m2 --- numerals}
Use numerals rather than words for quantities throughout (``sixteen of eighteen'',
``seventy-six seconds'', ``six hundred'').
\end{point}

\begin{reply}
Done throughout, by sweep and manual check: ``$16$ of $18$'', ``$600$ pairs'',
``$2{,}400$'', ``$2{,}000$ seeds'', ``$+0.42$, $+0.60$, $+1.14$, $+1.85$ nats at
$r = 1, 2, 4, 8$'', and so on. The seventy-six-second figure is withdrawn entirely
(R2.9). Number words are retained only where they are ordinary English rather than
reported quantities (``one of the two options'', ``three obstacles'').
\end{reply}

\begin{point}{R2--m3 --- narrative register and Section 1 length}
The register in Sections 1 and 5 is engaging, but Section 1 buries the contribution list
under three pages of motivation. Consider tightening.
\end{point}

\begin{reply}
The introduction is restructured with signposted paragraphs --- \emph{Three obstacles},
\emph{The design}, \emph{What the certificate claims, and what it does not},
\emph{Corrections carried by this version}, \emph{Contributions} --- so a reader reaches
the claims quickly and can navigate to the scope statement directly. The motivation is
compressed by roughly a third: the three-obstacle discussion is tightened and the
evaluation-critique literature now carries the motivational weight in two sentences with
citations rather than in expository prose. We have kept the narrative register in the
experiments, including the development anecdote about the NPO sign error, because it
carries a point about what statistical certification buys that we could not make as
economically any other way --- but we have removed the more self-conscious framings
(``One episode during development is worth recounting'' is now a plain topic sentence).
\end{reply}

\begin{point}{R2--m4 --- conservatism is not tightness}
A realised error of 0.031 against a nominal 0.05 is expected of a supermartingale and is
not itself evidence of tightness; the conservatism also costs power. Worth a sentence
acknowledging this.
\end{point}

\begin{reply}
Added, in \S6.2 immediately after the calibration numbers: ``A realised error of
$0.031$ against a nominal $0.05$ is what a supermartingale bound is \emph{supposed} to
produce; it is evidence of validity, not of tightness, and the same conservatism costs
power, which is the price Table~5 quantifies. We do not claim the bound is tight and the
gap is not a defect.'' The cross-reference is to the new sequential-baselines table, so
the reader can see the cost measured rather than merely conceded.
\end{reply}

\begin{point}{R2--m5 --- Fig.\ 2's left panel duplicates the text}
Fig.\ 2's left panel and the corresponding text make the same point twice; one could go.
\end{point}

\begin{reply}
The soundness result is now carried by the table and Remark~4 alone, and the redundant
figure panel is dropped from the manuscript.
\end{reply}

\begin{point}{R2--m6 --- Theorem 3's statement}
Theorem 3's statement should specify that $\mathbb{E}[\tau^*]$ is under exact unlearning
and small $\eps$, as the proof's final step assumes.
\end{point}

\begin{reply}
Done: Theorem~4 now reads ``\ldots the expectation being taken under exact unlearning and
the asymptotics in small $\eps$'', and the statement additionally notes that for the
finite-cohort target the time is capped by the cohort size $m$, since the claim resolves
deterministically once the cohort is exhausted.
\end{reply}

\begin{point}{R2--m7 --- typo in Section 5.3}
``$\Delta = 0.1$ in six hundred and ninety'' --- please check against Fig.\ 3.
\end{point}

\begin{reply}
Checked against the released simulation output: the value is $690$ pairs, which matches
Figure~3. The sentence now reads ``$\Delta=0.1$ in $690$'', in numerals.
\end{reply}

\clearpage
\section*{A note on what this revision does not settle}

\noindent Three of the reviewers' requests we have answered only partially, and we
would rather list them than let them be inferred from the text.

\begin{enumerate}[leftmargin=1.4em,itemsep=3pt]
\item \textbf{A LiRA-style attack outside $\F$} (R2.W4). Not run; it needs dozens of
retrained shadow references per cell. The two out-of-class attacks we did run are
weaker, and one of them already exceeds the tolerance on $3.2\%$ of certified runs, so
we regard the class restriction as demonstrated to be load-bearing but not fully
mapped.
\item \textbf{A paraphrase-aware score in $\F$} (R1.6, prompted by
\citealt{li2026beliefs}). Not implemented. It is the single most valuable extension we
can identify, precisely because the squeezing effect describes leakage our current class
cannot see, and we say so in both the related work and the conclusion.
\item \textbf{Method and probe coverage} (R1.5b, R1.5d). SimNPO is added; RMU and
task-vector negation are not, and the R-VOUCH probes remain anchored on NPO in the
benchmark tier. A general capability probe is added but it is \texttt{world\_facts} and
the MUSE holdout, not a full suite such as MMLU.
\end{enumerate}

\noindent Set against those, the changes we regard as materially strengthening the
paper are the ones the reviewers forced: a certificate whose validity no longer depends
on an independence assumption that our own data violates, a tolerance sweep that reaches
$\eps=0.05$ on real models with the cohort/super-population distinction made explicit, a
comparison against valid sequential procedures rather than a straw baseline, a
head-to-head against the metrics we claim to replace, an out-of-class attack that finds
a real gap, and a detectability measurement that overturns a claim we should never have
made without testing. We are grateful for both reports.

\bibliographystyle{plainnat}
\bibliography{refs}

\end{document}
